% !PRINTED 127
\weberpage{127}
%<<CONT
eine Potenz von $p$ ist. $R'$ ist der Inbegriff aller mit $Q'$ vertauschbaren Elemente von $P$.

Wir wenden nun weiter den Satz 3., \S.\,28 an, indem wir unter $P$, $R$, $Q$ die Gruppen verstehen, die hier mit denselben Buchstaben bezeichnet sind. Da hier $Q$ ein Theiler von $R$ ist, so ist $h_1 = 1$. Von den conjugirten Gruppen $b^{-1}Rb$ ist aber keine ausser $R$ durch $Q$ theilbar, und also sind die übrigen $h_2, h_3, \ldots$ alle grösser als $1$. Sie sind aber, da sie Theiler des Grades von $Q$ sein müssen, Potenzen von $p$ und aus der Gleichung

\[
j = h_1 + h_2 + h_3 + \cdots
\]

folgt also $j \equiv 1\ (\mathrm{mod}\ p)$.

Ist aber $Q_1$ irgend ein Divisor von $P$, dessen Grad eine Potenz von $p$ ist, so wenden wir wieder die Zerlegung von $j$ nach dem Theorem 3., \S.\,28 an, indem wir für die dort vorkommenden Gruppen $P$, $R$, $Q$ die Gruppen $P$, $R$, $Q_1$ setzen. Dann ist wieder

\[
j = h_1 + h_2 + h_3 + \cdots,
\]

worin jetzt $h_1, h_2, h_3, \ldots$ die Indices (in Bezug auf $Q$) der Durchschnitte von $Q_1$ mit der Gruppe $R$ und seinen conjugirten bedeuten, die also auch Potenzen von $p$ sind. % [sic: Weber prints "in Bezug auf Q"; after substituting Q_1 for theorem 3's Q, these indices are with respect to Q_1.]
Da aber $j \equiv 1\ (\mathrm{mod}\ p)$ ist, so muss mindestens eine von den Zahlen $h = 1$ sein, d. h. eine der Gruppen $b_\lambda^{-1}Rb_\lambda$ und folglich auch eine der Gruppen (11) ist durch $Q_1$ theilbar. Ist $Q_1$ vom Grade $p^\varkappa$, so ist es mit einer der Gruppen (11) identisch.

Damit ist also das Theorem II vollständig bewiesen.

\begin{center}
\S.\,31.

\gesperrt{Gruppen vom Grade $p^\alpha$.}
\end{center}

\gesperrt{Sylow} hat aus seinen Sätzen eine merkwürdige Eigenschaft der Gruppen abgeleitet, deren Grad eine Potenz einer Primzahl $p$, $p^\alpha$, ist. Es sei also $P$ eine solche Gruppe und $\beta$ eine positive Zahl, kleiner als $\alpha$. Dann enthält nach dem Satze I, \S.\,29 die Gruppe $P$ einen Theiler $Q$ vom Grade $p^\beta$, deren Index $p^{\alpha-\beta}$ ist. Wir wenden den Satz 3., \S.\,28 an, indem wir für $R$ und $Q$ diese Gruppe $Q$ vom Grade $p^\beta$ setzen, dann sind $h_1, h_2, h_3, \ldots$
%CONT>>
